% Lösungen Einheit 1 von 3 – Ansatz und Punktbedingungen (zusammenbau v0.1)
\einheitenkopf{Einheit 1 von 3 · Ansatz und Punktbedingungen}

\erg{9}{a) genau passend – drei Unbekannte $a$, $b$, $c$ und drei Gleichungen. \quad b) $m = 3$, $n = 2$, also $y = 3x + 2$ \quad c) $h(0{,}6) = 1{,}2$ und $h(3) = 6$ – beide Punkte liegen auf dem Graphen, und durch zwei Punkte geht genau eine Gerade. \quad d) $c = 2$; $a + b = -3$ und $9a + 3b = 3$; $a = 2$, $b = -5$, also $g(x) = 2x^2 - 5x + 2$ \quad e) Ansatz $g(x) = ax^2 + c$: $a + c = -3$ und $9a + c = 13$; $a = 2$, $c = -5$, also $g(x) = 2x^2 - 5$ \quad f) Nullstellen $-2$, $1$ und $3$: $f(x) = a(x + 2)(x - 1)(x - 3)$; $f(0) = 6a = 3$, also $a = 0{,}5$ und $f(x) = 0{,}5(x + 2)(x - 1)(x - 3)$ \quad g) $f(0) = a = 2$; $f(1) = 2b = 6$, also $b = 3$ \quad h) $(-1|0)$, $(1|0)$, $(0|1{,}8)$; $g(x) = ax^2 + 1{,}8$ mit $g(1) = 0$: $a = -1{,}8$, also $g(x) = -1{,}8x^2 + 1{,}8$ \quad i) Ansatz $f(x) = ax^4 + bx^2 + c$, $f'(x) = 4ax^3 + 2bx$: $c = -2$; $f(2) = 16a + 4b - 2 = -6$ und $f'(2) = 32a + 4b = 0$; $a = 0{,}25$, $b = -2$, also $f(x) = 0{,}25x^4 - 2x^2 - 2$}
\erg{10}{a) Die Symmetrie zur $y$-Achse streicht $b$; richtig: $f(x) = ax^2 + c$ mit $a + c = 4$ und $4a + c = 13$, also $a = 3$, $c = 1$ und $f(x) = 3x^2 + 1$. \quad b) Achsensymmetrie heißt $f(-x) = f(x)$; Glieder mit ungeradem Exponenten wechseln beim Einsetzen von $-x$ das Vorzeichen, darum dürfen nur gerade Exponenten vorkommen.}
